% ============================================================
%  English version — main document
%  Universidad de Buenos Aires · FCEN · Departamento de Computación
%  Licenciatura en Ciencias de la Computación
% ============================================================
\documentclass[11pt,a4paper,twoside]{tesis}

\usepackage{graphicx}
\usepackage[utf8]{inputenc}
\usepackage[T1]{fontenc}
\usepackage[english]{babel}
\usepackage[left=3cm,right=3cm,bottom=3.5cm,top=3.5cm]{geometry}
\usepackage{amsmath,amssymb,amsthm}
\usepackage{hyperref}
\usepackage{booktabs}
\usepackage{multirow}
\usepackage{float}
\usepackage{algorithm}
\usepackage{algpseudocode}
\usepackage[backend=biber,style=numeric,sorting=nyt]{biblatex}

\addbibresource{tesis.bib}

\hypersetup{colorlinks=true,linkcolor=black,citecolor=black,urlcolor=black}

% Redefine chapter/section names for English
\def\contentsname{Contents}
\def\listfigurename{List of Figures}
\def\listtablename{List of Tables}
\def\bibname{Bibliography}
\def\indexname{Index}
\def\figurename{Figure}
\def\tablename{Table}
\def\chaptername{Chapter}
\def\appendixname{Appendix}
\def\abstractname{Abstract}

% Theorem environments
\newtheorem{definition}{Definition}[chapter]
\newtheorem{remark}{Remark}[chapter]

% Notation shorthands
\newcommand{\G}{\mathcal{G}}
\newcommand{\V}{\mathcal{V}}
\newcommand{\E}{\mathcal{E}}
\newcommand{\C}{\mathcal{C}}
\newcommand{\Gnew}{\mathcal{G}'}
\newcommand{\Cnew}{\mathcal{C}'}
\newcommand{\simf}{\mathrm{sim}}
\newcommand{\SR}{\mathrm{SR}}
\newcommand{\NMI}{\mathrm{NMI}}

\begin{document}

%%%% COVER PAGE
\def\autor{Dante Culaciati}
\def\tituloTesis{Improving Community Membership Hiding \vspace{.2cm} \\
                 via Deep Reinforcement Learning}
\def\runtitulo{Mejorando el Ocultamiento de Membres\'ia de Comunidades
               mediante Aprendizaje Profundo por Refuerzo}
\def\runtitle{Improving Community Membership Hiding via Deep Reinforcement Learning}
\def\director{Esteban Feuerstein}
\def\codirector{[Nombre del Codirector]}   % TODO: fill in
\def\lugar{Buenos Aires, 2026}
\input{caratula}

%%%% FRONT MATTER
\frontmatter
\pagestyle{empty}
\chapter*{\runtitulo}

\noindent
El problema de \emph{Community Membership Hiding} (CMH) consiste en, dado un grafo, un
algoritmo de detecci\'on de comunidades y un nodo objetivo $u$, sugerir un conjunto m\'inimo de
reconexiones de aristas para que $u$ deje de ser detectado como miembro de su comunidad
original. Este problema es relevante en el contexto de privacidad en redes, donde los
algoritmos de detecci\'on pueden revelar informaci\'on sensible sobre los usuarios a partir
de su posici\'on estructural en el grafo.

Trabajos recientes han formulado CMH como un Proceso de Decisi\'on de Markov resuelto
mediante Aprendizaje Profundo por Refuerzo (DRL), utilizando una arquitectura Actor-Cr\'itico.
En este trabajo identificamos una limitaci\'on arquitect\'onica cr\'itica del m\'etodo original: el
agente eval\'ua el grafo de forma global sin condicionar expl\'icitamente en el nodo objetivo,
lo que le impide aprender una pol\'itica espec\'ifica de ocultamiento.

Proponemos una arquitectura DRL novedosa que incorpora expl\'icitamente la identidad del
nodo objetivo como entrada de la red, utiliza \emph{Graph Attention Networks v2} (GATv2)
como codificador de grafos, y combina se\~nales est\'aticas y din\'amicas para representar
fielmente el estado del grafo durante el proceso de reconexci\'on. Extendemos este sistema
de un solo agente a una \emph{Arquitectura Dual}, que especializa agentes separados para la
adici\'on y eliminaci\'on de aristas, e introducimos una heur\'istica de \emph{Subgrafado Guiado
por Pol\'itica} que hace escalable el m\'etodo en grafos de gran tama\~no.

La evaluaci\'on experimental sobre los datasets est\'andar del \'area muestra que nuestra
arquitectura supera consistentemente al m\'etodo original y a los m\'etodos de referencia
basados en gradientes, logrando altas tasas de \'exito de ocultamiento con una perturbaci\'on
estructural m\'inima.

\bigskip

\noindent\textbf{Palabras clave:} detecci\'on de comunidades, ocultamiento de comunidades,
aprendizaje por refuerzo profundo, redes de atenci\'on de grafos, privacidad en redes,
GATv2, Actor-Cr\'itico.


\cleardoublepage
\chapter*{\runtitle}

\noindent
The \emph{Community Membership Hiding} (CMH) problem asks: given a graph, a community
detection algorithm, and a target node $u$, find a minimal set of edge rewirings such
that $u$ is no longer detected as a member of its original community. This problem arises
naturally in the context of network privacy, where community detection algorithms can
inadvertently expose sensitive user affiliations---such as political or ideological
leanings---solely from the structural position of a node in the underlying graph.

Recent work has formulated CMH as a Markov Decision Process (MDP) solved via Deep
Reinforcement Learning (DRL), employing an Actor-Critic architecture. While that
approach demonstrated the viability of the paradigm, we identify a critical architectural
limitation: the agent evaluates the network globally without explicitly conditioning on
the identity of the target node, thereby learning a generic community-shuffling policy
rather than a targeted hiding strategy.

In this thesis, we propose a novel DRL architecture that addresses these fundamental
shortcomings. Our model explicitly incorporates the target node as a feature of the
network representation, employs Graph Attention Networks v2 (GATv2) as the graph
encoder to leverage dynamic, query-conditioned attention, and combines static and
real-time structural signals to accurately track topological changes during rewiring.
We further extend the single-agent system into a \emph{Dual-Agent Architecture} that
dedicates separate agents to edge addition and edge deletion, and introduce a
\emph{Policy-Guided Subgraphing} heuristic that makes the framework tractable on
large-scale graphs.

Extensive experimental evaluation on the standard benchmark datasets of the field
demonstrates that our architecture consistently outperforms both the original DRL
agent and strong gradient-based baselines, achieving high hiding success rates with
minimal structural perturbation.

\bigskip

\noindent\textbf{Keywords:} community detection, community membership hiding, deep
reinforcement learning, graph attention networks, network privacy, GATv2, Actor-Critic.


\cleardoublepage
\chapter*{Acknowledgements}

% TODO: complete acknowledgements
\noindent
[To be completed.]


\cleardoublepage
\input{dedicatoria.tex}

\cleardoublepage
\tableofcontents

% ============================================================
\mainmatter
\pagestyle{headings}
% ============================================================


% ============================================================
\chapter{Introduction}
\label{ch:intro}
% ============================================================

The modern science of networks has brought significant advances to our understanding of
complex systems. One of the most relevant structural features of graphs representing
real-world phenomena---ranging from social media platforms to biological interaction
networks---is their \emph{community structure}: the tendency for nodes to cluster into
densely connected groups that are only sparsely connected to the rest of the network.
Community detection algorithms that uncover this structure have proven highly effective
\cite{Fortunato-2010}, and are now routinely deployed in large-scale platforms to
characterize users, personalize content, and target advertisements.

However, this functionality carries a fundamental privacy risk. When a graph encodes
social relationships, the community a node belongs to may inadvertently reveal sensitive
affiliations---political leanings, religious beliefs, health conditions, or ideological
associations---purely from its structural position, without any explicit disclosure of
such information. This concern is not merely theoretical: community detection has been
used in practice to identify political groups in online networks and to profile individuals
in social graphs.

To counter this, the concept of \emph{Community Membership Hiding} (CMH) was introduced
by \textcite{bernini-2024}. The goal of CMH is to strategically modify the topology of a
network---by adding or removing a small number of edges adjacent to a target node $u$---so
that $u$ is no longer assigned to its original community by a detection algorithm, while
keeping overall graph disruption minimal. The detection algorithm is treated as a
\emph{black box}, meaning that no knowledge of its internals is exploited.

Bernini et al.\ formulated CMH as a Markov Decision Process (MDP) and solved it via Deep
Reinforcement Learning (DRL) using an Actor-Critic architecture (A2C). While that work
demonstrated the viability of the DRL approach, it exhibits notable performance degradation
on larger graphs, and---as we show in this thesis---harbors a critical architectural
flaw: the agent evaluates the network \emph{globally}, without explicitly conditioning
on the identity of the target node. As a consequence, the agent cannot learn a policy
that depends on who the target is; it instead learns a generic community-shuffling
strategy that happens to displace the target node as a side effect.

\section*{Contributions}

In this thesis, we make the following contributions:

\begin{enumerate}
    \item \textbf{Identification of an architectural limitation.}
          We formally demonstrate that the original DRL agent produces identical outputs
          for different target nodes given the same graph state, confirming that it cannot
          condition its policy on the target.

    \item \textbf{Target-aware architecture.}
          We propose a novel DRL architecture that explicitly encodes the identity of the
          target node as an additional input feature. The graph encoder is replaced by
          Graph Attention Networks v2 (GATv2) \cite{brody2022how}, which provide a more
          expressive, query-conditioned attention mechanism compared to the original GNN.
          Node features combine static indicators derived from the original graph with
          dynamically updated structural signals that reflect topological changes as
          rewiring proceeds.

    \item \textbf{Dual-Agent Architecture.}
          We extend the single-agent system into a Dual-Agent Architecture that maintains
          two specialized agents---one for edge addition and one for edge deletion---with
          separate policy networks and learning dynamics. We show that this specialization
          leads to improved performance and learning stability.

    \item \textbf{Policy-Guided Subgraphing.}
          To make the framework tractable on large graphs, we introduce a
          Policy-Guided Subgraphing heuristic that restricts the action space to a
          relevant neighborhood around the target node, reducing computational cost while
          preserving the quality of the learned policy.

    \item \textbf{Extensive experimental evaluation.}
          We evaluate all proposed contributions on the standard benchmark datasets used
          in the CMH literature, comparing against the original DRL agent and two
          gradient-based methods introduced in \textcite{silvestri2025righthidemaskingcommunity}.
\end{enumerate}

\section*{Organization}

The remainder of this thesis is organized as follows. Chapter~\ref{ch:background}
provides the necessary background on graphs, community detection, reinforcement learning,
and graph neural networks. Chapter~\ref{ch:problem} gives a formal problem statement.
Chapter~\ref{ch:related} reviews related work.
Chapter~\ref{ch:architecture} describes the proposed architecture. Chapter~\ref{ch:experiments}
presents the experimental evaluation. Chapter~\ref{ch:conclusions} concludes with a
discussion and directions for future work.


% ============================================================
\chapter{Background}
\label{ch:background}
% ============================================================

This chapter introduces the key concepts required to understand the rest of the thesis.
We cover the fundamentals of graphs and community detection, provide a formal definition
of Community Membership Hiding, review the relevant machinery of Reinforcement Learning,
and describe Graph Attention Networks and the node2vec algorithm.

% ------------------------------------------------------------
\section{Graphs and Community Structure}
\label{sec:graphs}
% ------------------------------------------------------------

A \emph{graph} $\G = (\V, \E)$ consists of a set of nodes (vertices) $\V$ and a set
of edges $\E \subseteq \V \times \V$. We denote $|\V| = N$ and $|\E| = M$. In an
undirected graph, $(u,v) \in \E$ implies $(v,u) \in \E$; in a directed graph, these are
distinct edges. Most of the datasets in this work are undirected.

The \emph{degree} of a node $v$, denoted $\deg(v)$, is the number of edges incident to
$v$. The mean degree of the graph is $\mu = 2M/N$. The \emph{neighborhood} of $v$ is
$\mathcal{N}(v) = \{w : (v,w) \in \E\}$.

\begin{definition}[Community structure]
A graph $\G$ exhibits \emph{community structure} if its nodes can be partitioned into
groups $\C = \{\C_1, \C_2, \ldots, \C_k\}$ such that the density of edges within each
group is significantly higher than the density of edges between different groups.
\end{definition}

Intuitively, communities correspond to densely connected subgraphs (clusters) that are
only loosely coupled to one another. This definition is intentionally informal; different
community detection algorithms operationalize it differently.

% ------------------------------------------------------------
\section{Community Detection Algorithms}
\label{sec:community-detection}
% ------------------------------------------------------------

A large number of algorithms for community detection have been proposed in the literature
\cite{Fortunato-2010}. In this thesis, we focus on three widely used methods, all
of which are treated as \emph{black boxes} by our agent.

\paragraph{Greedy modularity optimization.}
\textcite{clauset2004greedy} introduced a greedy algorithm that maximizes \emph{modularity},
defined as the fraction of edges within communities minus the expected fraction under a
null model with the same degree sequence:
\[
Q = \frac{1}{2M}\sum_{i,j}\left[A_{ij} - \frac{k_i k_j}{2M}\right]\delta(c_i, c_j),
\]
where $A$ is the adjacency matrix, $k_i$ is the degree of node $i$, and $\delta(c_i,c_j)$
is 1 if nodes $i$ and $j$ belong to the same community and 0 otherwise. The algorithm
begins with each node in its own community and iteratively merges pairs of communities
that yield the largest increase in $Q$.

\paragraph{Louvain algorithm.}
\textcite{blondel2008louvain} proposed a two-phase greedy heuristic for modularity
maximization. In the first phase, each node is moved to the community of a neighbor that
yields the largest modularity gain (if any such gain exists). In the second phase,
communities are collapsed into super-nodes and the first phase is repeated on the
resulting graph. The process continues until no further improvement is possible. The
Louvain algorithm is highly efficient and widely adopted for large-scale networks.

\paragraph{Walktrap algorithm.}
\textcite{pons2005walktrap} introduced a community detection algorithm based on random
walks. The key intuition is that short random walks tend to stay within the same community;
therefore, nodes that are visited by the same set of short random walks are likely to
belong to the same community. The algorithm defines a distance between nodes based on the
transition probabilities of a $t$-step random walk and uses hierarchical clustering to
identify communities.

% ------------------------------------------------------------
\section{Community Membership Hiding: Informal Overview}
\label{sec:cmh-informal}
% ------------------------------------------------------------

In the Community Membership Hiding setting, an actor (the \emph{hider}) controls a
specific target node $u$ in a graph $\G$. The hider can add or remove a limited number
of edges adjacent to $u$, with the goal of causing a given community detection algorithm
$f$ to no longer assign $u$ to its original community $\C_i$ after the modifications.
The detection algorithm is treated as a black box: the hider can query it but cannot
access its internals.

This problem models a privacy-conscious user who wants to prevent a platform from
correctly inferring their social group, and who can strategically manage their
connections to achieve this. A formal treatment is given in Chapter~\ref{ch:problem}.

% ------------------------------------------------------------
\section{Reinforcement Learning}
\label{sec:rl}
% ------------------------------------------------------------

\emph{Reinforcement Learning} (RL) is a machine learning paradigm in which an
\emph{agent} learns to make sequential decisions by interacting with an \emph{environment}
in order to maximize a cumulative \emph{reward} signal. The interaction proceeds in
discrete time steps: at each step $t$, the agent observes a state $s_t \in \mathcal{S}$,
selects an action $a_t \in \mathcal{A}(s_t)$ according to a policy $\pi$, and receives
a scalar reward $r_t$ while transitioning to state $s_{t+1}$.

\paragraph{Markov Decision Process.}
The standard formalization is the Markov Decision Process (MDP), defined by the tuple
$(\mathcal{S}, \mathcal{A}, P, R, \gamma)$, where $\mathcal{S}$ is the state space,
$\mathcal{A}$ is the action space, $P(s' \mid s, a)$ is the transition probability
distribution, $R(s, a)$ is the expected immediate reward, and $\gamma \in [0,1]$ is the
discount factor. The return from step $t$ is:
\[
G_t = \sum_{k=0}^{\infty} \gamma^k r_{t+k+1}.
\]
The objective is to find a policy $\pi^*$ that maximizes the expected return:
\[
\pi^* = \arg\max_{\pi} \; \mathbb{E}_{\pi}\left[G_t \mid s_t = s\right].
\]

\paragraph{Actor-Critic methods.}
Policy gradient methods optimize a parametric policy $\pi_\theta(a \mid s)$ by following
the gradient of the expected return:
\[
\nabla_\theta J(\theta) = \mathbb{E}_{s,a}\left[\nabla_\theta \log \pi_\theta(a \mid s)
    \, A^\pi(s, a)\right],
\]
where $A^\pi(s, a) = Q^\pi(s, a) - V^\pi(s)$ is the \emph{advantage function}, which
measures how much better action $a$ is compared to the average action in state $s$.

Actor-Critic architectures maintain two networks: the \emph{actor}, which represents the
policy $\pi_\theta$, and the \emph{critic}, which estimates the state-value function
$V_\phi(s)$ used to compute the advantage. Training the critic to provide an accurate
baseline substantially reduces the variance of the policy gradient estimate without
introducing bias, leading to more stable learning.

The Advantage Actor-Critic (A2C) algorithm \cite{mnih2016asynchronous}, which forms the
basis of the original CMH agent and of our proposed architecture, synchronously updates
the actor and critic using the advantage estimated from collected trajectories.

% ------------------------------------------------------------
\section{Graph Neural Networks}
\label{sec:gnn}
% ------------------------------------------------------------

Graph Neural Networks (GNNs) are a family of deep learning models designed to operate
on graph-structured data. The core idea is \emph{message passing}: each node aggregates
information from its neighbors, transforming the aggregated messages through learnable
parameters to produce a new node representation. Formally, the $l$-th layer of a GNN
updates the representation of node $v$ as:
\[
\mathbf{h}_v^{(l)} = \mathrm{UPDATE}^{(l)}\!\left(\mathbf{h}_v^{(l-1)},\,
    \mathrm{AGG}^{(l)}\!\left(\left\{\mathbf{h}_u^{(l-1)} : u \in \mathcal{N}(v)\right\}\right)\right),
\]
where the initial representation $\mathbf{h}_v^{(0)}$ is the feature vector of node $v$.

\subsection{Graph Attention Networks (GAT)}

Graph Attention Networks \cite{veličković2018graph} replace the uniform or degree-based
aggregation of vanilla GNNs with a learnable attention mechanism. The $l$-th layer
computes:
\[
\mathbf{h}_v^{(l)} = \sigma\!\left(\sum_{u \in \mathcal{N}(v)} \alpha_{vu}\,
    \mathbf{W}^{(l)} \mathbf{h}_u^{(l-1)}\right),
\]
where the attention coefficient $\alpha_{vu}$ is computed by a small neural network
applied to the concatenation of the projected features of $v$ and $u$:
\[
\alpha_{vu} = \frac{\exp\!\left(\mathrm{LeakyReLU}\!\left(\mathbf{a}^\top
    [\mathbf{W}\mathbf{h}_v \,\|\, \mathbf{W}\mathbf{h}_u]\right)\right)}
    {\sum_{w \in \mathcal{N}(v)}\exp\!\left(\mathrm{LeakyReLU}\!\left(\mathbf{a}^\top
    [\mathbf{W}\mathbf{h}_v \,\|\, \mathbf{W}\mathbf{h}_w]\right)\right)}.
\]

\subsection{GATv2: Dynamic Graph Attention}

\textcite{brody2022how} identified a fundamental limitation of the original GAT: its
attention is \emph{static}, meaning the ranking of neighbors is the same for all query
nodes. This occurs because the attention scoring function applies the linear transformation
\emph{before} the non-linearity, making the relative importance of neighbors independent
of the query node's feature.

GATv2 corrects this by reordering the operations:
\[
e(\mathbf{h}_v, \mathbf{h}_u) = \mathbf{a}^\top \mathrm{LeakyReLU}\!\left(
    \mathbf{W}_1 \mathbf{h}_v + \mathbf{W}_2 \mathbf{h}_u\right),
\]
which makes the scoring function a two-layer MLP over the concatenated features. The
resulting attention is \emph{dynamic}: the ranking of neighbors can differ for every
query node, making GATv2 a strictly more expressive model. This property is particularly
desirable in our setting, where the importance of each neighbor to the target node should
depend on the target's own features.

% ------------------------------------------------------------
\section{The node2vec Algorithm}
\label{sec:node2vec}
% ------------------------------------------------------------

In many real-world graphs, nodes have no natural feature vectors. node2vec
\cite{grover2016node2vec} addresses this by learning a $d$-dimensional embedding for each
node via a procedure analogous to word2vec. The algorithm generates sequences of nodes
by performing \emph{biased second-order random walks} on the graph, then trains a
skip-gram model with negative sampling on these sequences.

The bias is controlled by two parameters:
\begin{itemize}
    \item $p$ (return parameter): probability of returning to the previously visited node.
      Low $p$ encourages a depth-first search behavior.
    \item $q$ (in-out parameter): probability of exploring nodes far from the previously
      visited node. Low $q$ encourages a breadth-first search behavior.
\end{itemize}
By tuning $p$ and $q$, node2vec can interpolate between capturing \emph{homophily}
(nodes in the same community are embedded close together) and
\emph{structural equivalence} (nodes playing similar structural roles are embedded
close together).

In the original CMH architecture, node2vec is used to generate initial node features.
However, these embeddings are computed once on the original graph and are never updated
as the graph changes during the rewiring process. This is one of the limitations we
address in our proposed architecture.


% ============================================================
\chapter{Problem Formulation}
\label{ch:problem}
% ============================================================

In this chapter we give a formal definition of the Community Membership Hiding problem,
specifying the constraints, the objective, the reward structure, and the evaluation
metrics.

% ------------------------------------------------------------
\section{Formal Definition}
\label{sec:formal}
% ------------------------------------------------------------

Let $\G = (\V, \E)$ be an undirected, unweighted graph and let $f$ be a community
detection algorithm that, given $\G$, produces a partition $\C = f(\G) =
\{\C_1, C_2, \ldots, C_k\}$ of $\V$ into $k$ non-overlapping communities.

Given a target node $u \in \V$ with original community assignment $\C_i = f_u(\G)$
(the community to which $u$ belongs in $\G$), the goal of CMH is to find a modified
graph $\Gnew = (\V, \E')$ such that the following \emph{deception condition} is satisfied:

\begin{definition}[Deception condition]
\label{def:deception}
Let $\C_i' = f_u(\Gnew)$ denote the community assigned to $u$ in $\Gnew$. The deception
condition is:
\[
\simf(\C_i \setminus \{u\},\; \C_i' \setminus \{u\}) \leq \tau,
\]
where $\simf$ is a community similarity function and $\tau \in [0,1]$ is the
\emph{deception threshold}. A value of $\tau = 0$ requires $\C_i$ and $\C_i'$ to be
completely disjoint (hard deception), while $\tau = 1$ permits them to be identical
(trivially satisfied).
\end{definition}

Following \textcite{bernini-2024}, we use the \emph{S{\o}rensen-Dice Coefficient} (DSC)
as $\simf$:
\[
\mathrm{DSC}(A, B) = \frac{2|A \cap B|}{|A| + |B|}.
\]
DSC measures the proportion of shared nodes between two communities: it equals 0 when
the communities are disjoint and 1 when they are identical.

\begin{definition}[Budget constraint]
\label{def:budget}
The number of edge modifications is bounded by a \emph{budget}
$\beta = \lfloor \beta_{\mathrm{mul}} \cdot \mu \rfloor$, where $\mu = 2M/N$ is the
mean degree and $\beta_{\mathrm{mul}} \in \{1/2, 1, 2\}$ is a multiplier. Each
rewiring step constitutes a single modification.
\end{definition}

\begin{definition}[Admissible actions]
\label{def:actions}
To align with the hiding objective, the action space is restricted as follows. Let
$\C_i = f_u(\G)$ be the current community of $u$:
\begin{itemize}
    \item \textbf{Edge removal}: the agent may remove an edge $(u, v)$ only if
      $v \in \C_i$ (i.e., $v$ is an \emph{intra-community} neighbor of $u$).
    \item \textbf{Edge addition}: the agent may add an edge $(u, v)$ only if
      $v \notin \C_i$ (i.e., $v$ is an \emph{inter-community} candidate).
\end{itemize}
Removing inter-community edges would further isolate $u$ within $\C_i$; adding
intra-community edges would reinforce $u$'s membership. Both are thus excluded.
\end{definition}

\section{Graph Similarity}

Beyond the community-level deception condition, we also measure the global structural
distortion introduced by the rewiring. Following the original paper, we use the
\emph{Jaccard similarity} of the edge sets:
\[
\mathrm{JAC}(\G, \Gnew) = \frac{|\E \cap \E'|}{|\E \cup \E'|}.
\]
A high Jaccard similarity indicates that the modified graph closely resembles the
original, i.e., the structural cost of hiding is low.

\section{Normalized Mutual Information}

To quantify the impact of the rewiring on the \emph{community structure} as a whole,
we use the Normalized Mutual Information (NMI) \cite{strehl2003cluster} between the
original partition $\C = f(\G)$ and the modified partition $\C' = f(\Gnew)$:
\[
\NMI(\C, \C') = \frac{2\, I(\C\,;\, \C')}{H(\C) + H(\C')},
\]
where $I$ denotes mutual information and $H$ entropy over the community assignments.
$\NMI$ takes values in $[0,1]$, with 1 indicating identical partitions and 0 indicating
no shared information. A high NMI after hiding means the global community structure has
been well preserved.

\section{Evaluation Metrics}
\label{sec:metrics}

We evaluate methods using three metrics:

\paragraph{Success Rate (SR).} The fraction of test instances in which the deception
condition (Definition~\ref{def:deception}) is satisfied:
\[
\SR = \frac{\#\,\text{successful runs}}{\#\,\text{total runs}}.
\]

\paragraph{NMI.} The average NMI between $f(\G)$ and $f(\Gnew)$ across all test
instances. Higher NMI indicates less structural distortion.

\paragraph{F1 score.} Since SR and NMI are competing objectives---higher SR typically
requires more aggressive rewiring, which lowers NMI---we report their harmonic mean as a
balanced metric:
\[
F_1 = \frac{2 \cdot \SR \cdot \NMI}{\SR + \NMI}.
\]


% ============================================================
\chapter{Related Work}
\label{ch:related}
% ============================================================

\section{The Original CMH Agent (Bernini et al., 2024)}
\label{sec:bernini}

\textcite{bernini-2024} were the first to formulate Community Membership Hiding as an
MDP and solve it with DRL. Their architecture, which serves as our primary baseline,
consists of:
\begin{enumerate}
    \item A \textbf{graph encoder} based on a GNN that produces a representation of the
          graph state. Node features are either random vectors or node2vec embeddings
          computed on the original graph.
    \item An \textbf{Actor-Critic network} (A2C) that, given the encoded graph, outputs
          a probability distribution over candidate nodes for the next rewiring step and a
          scalar value estimate.
\end{enumerate}
The agent is trained on the \texttt{words} dataset using the \texttt{greedy} community
detection algorithm, and is evaluated for transferability to other datasets and algorithms.

Notably, the original agent was shown to achieve success rates significantly above the
random and degree-based baselines in standard evaluation scenarios. However, it degrades
substantially on larger graphs, and---as we establish in Chapter~\ref{ch:architecture}---
it suffers from a fundamental architectural limitation regarding the role of the target
node.

\section{The Right to Hide (Silvestri et al., 2025)}
\label{sec:nabla}

\textcite{silvestri2025righthidemaskingcommunity} approached CMH from a different angle,
framing it as a continuous optimization problem. They define a differentiable relaxation
of the deception loss and apply projected gradient descent, yielding two variants:
$\nabla$-CMH and $\nabla$-CMH (projected). These gradient-based methods are
computationally efficient and achieve competitive hiding rates with a well-controlled
structural footprint. We use both as strong baselines in our evaluation.

\section{Evasion Problems on Graphs}

CMH is part of a broader family of \emph{graph evasion problems}, in which a structural
modification is sought to fool an algorithm operating on the graph. Related problems
include adversarial attacks on node classification \cite{kipf2017semi}, link prediction
evasion, and graph-level classification adversarial examples. Our work is specifically
concerned with the community-detection variant and does not assume any knowledge of the
internal parameters of the detection algorithm.


% ============================================================
\chapter{Proposed Architecture}
\label{ch:architecture}
% ============================================================

In this chapter we first analyze the original DRL agent and identify its key limitations.
We then describe our proposed architecture step by step.

% ------------------------------------------------------------
\section{Analysis of the Original Agent}
\label{sec:analysis}
% ------------------------------------------------------------

\subsection{The Target-Node Blindness Problem}

A central observation motivating our work is the following. In the original architecture
\cite{bernini-2024}, the graph encoder receives the full graph state $\G_t$ as input
and produces a vector representation for every node. The actor network then outputs a
probability distribution over all candidate nodes, and the environment applies the
selected action \emph{relative to the target node $u$} (e.g., removes the edge between
$u$ and the selected intra-community neighbor).

The critical issue is that nowhere in this pipeline is the identity of $u$ communicated
to the network. The encoder and actor receive no signal about \emph{who} the target is.
To verify this, we fixed all random seeds and ran the original agent on the same graph
with different target nodes. The actor's output probabilities and the critic's value
estimate were identical across all target choices in the first step of every episode,
confirming that the model is entirely unaware of $u$.

As a result, the agent cannot adapt its policy to the specific structural position of the
target node. Instead, it learns a global rewiring strategy that tends to increase overall
community mixing---effectively shuffling the entire graph's community structure. While
this can accidentally displace $u$ from its community, it is neither efficient nor
reliable, especially when the budget $\beta$ is small.

\subsection{Static Node Embeddings}

The original architecture uses node2vec embeddings computed \emph{once} on the original
graph $\G$ and reused throughout the entire episode. As the rewiring process modifies
the graph, the embeddings become increasingly stale and fail to reflect the current
topology. Furthermore, computing node2vec requires a separate preprocessing step and is
itself computationally expensive for large graphs.

\subsection{Undifferentiated Action Space}

The original agent handles both edge addition and edge removal with the same policy
network, without distinguishing between the two types of actions. We hypothesize that
these two tasks have fundamentally different structure: removing intra-community edges
requires identifying which neighbors contribute most to $u$'s community membership,
while adding inter-community edges requires identifying which external nodes are most
useful for attracting $u$ to a different community. Conflating these two tasks may limit
the expressiveness of the learned policy.

% ------------------------------------------------------------
\section{Target-Aware Node Feature Construction}
\label{sec:features}
% ------------------------------------------------------------

Our first modification addresses the target-node blindness problem. For each node $v \in
\V$, we construct a feature vector $\mathbf{x}_v \in \mathbb{R}^d$ that explicitly
encodes whether $v$ is the target node:
\[
\mathbf{x}_v = \left[\,f_{\mathrm{struct}}(v) \;\Big|\; \mathbf{1}[v = u]\,\right],
\]
where $f_{\mathrm{struct}}(v)$ is a vector of structural features and $\mathbf{1}[v=u]$
is a binary indicator that is 1 only for the target node $u$.

The structural features $f_{\mathrm{struct}}(v)$ combine:
\begin{itemize}
    \item \textbf{Static features} computed on the original graph $\G_0$ (degree,
          clustering coefficient, betweenness centrality, community label) and held
          fixed throughout the episode. These capture the initial structural role of
          each node.
    \item \textbf{Dynamic features} recomputed at each step on the current graph $\G_t$
          (current degree, current community label, current number of intra-community
          and inter-community neighbors). These capture the real-time effect of previous
          rewiring actions.
\end{itemize}

By including $\mathbf{1}[v = u]$ in the feature vector, the GATv2 encoder can compute
attention weights that depend explicitly on whether a node is the target. This enables
the model to give special treatment to $u$'s neighborhood and to aggregate information
relevant to the hiding objective.

% ------------------------------------------------------------
\section{GATv2-Based Graph Encoder}
\label{sec:encoder}
% ------------------------------------------------------------

We replace the original GNN encoder with a multi-layer GATv2 \cite{brody2022how}
network. As discussed in Section~\ref{sec:gnn}, GATv2 provides dynamic attention:
the relative importance of each neighbor can differ for every query node, which is
essential when the target node $u$ needs to focus on a specific subset of its neighbors.

The GATv2 encoder takes as input the current graph $\G_t$ with node features
$\mathbf{X}_t \in \mathbb{R}^{N \times d}$ and produces a node embedding matrix
$\mathbf{H} \in \mathbb{R}^{N \times d'}$. The actor and critic networks operate on
these embeddings.

\paragraph{Actor.}
The actor scores each candidate node $v$ in the action set $\mathcal{A}(s_t)$:
\[
\mathrm{score}(v) = \mathrm{MLP}_{\mathrm{actor}}\!\left(\left[\mathbf{h}_u \,\|\, \mathbf{h}_v\right]\right),
\]
where $\mathbf{h}_u$ and $\mathbf{h}_v$ are the embeddings of the target node and
candidate $v$ respectively, and $\|\,\|$ denotes concatenation. This formulation ensures
that the score of each candidate is conditioned on the target node's embedding, directly
linking the action selection to the target's current context.

The action probability distribution is obtained by applying a softmax over all
candidates in $\mathcal{A}(s_t)$.

\paragraph{Critic.}
The critic estimates the state value:
\[
V(s_t) = \mathrm{MLP}_{\mathrm{critic}}\!\left(\mathrm{READOUT}(\mathbf{H})\right),
\]
where READOUT is a mean pooling over all node embeddings, followed by a linear
transformation.

% ------------------------------------------------------------
\section{Dual-Agent Architecture}
\label{sec:dual}
% ------------------------------------------------------------

The Dual-Agent Architecture maintains two independent Actor-Critic systems:
\begin{itemize}
    \item $\mathrm{Agent}_{\mathrm{add}}$: handles edge \emph{addition} actions
      (selecting inter-community candidates to connect to $u$).
    \item $\mathrm{Agent}_{\mathrm{del}}$: handles edge \emph{removal} actions
      (selecting intra-community neighbors of $u$ whose edge should be removed).
\end{itemize}

At each step, the environment determines whether the next action should be an addition
or a deletion (based on availability and the current budget), and dispatches to the
appropriate agent. Each agent has its own GATv2 encoder, actor network, and critic
network, with parameters trained independently.

The rationale is that edge addition and edge deletion are structurally different tasks.
The deletion agent needs to identify which intra-community connections most strongly
anchor $u$ to its community; the addition agent needs to identify which external nodes
would most effectively attract $u$ to a different community. Separating these tasks
into specialized networks allows each agent to develop representations tailored to
its subtask.

% ------------------------------------------------------------
\section{Policy-Guided Subgraphing}
\label{sec:subgraph}
% ------------------------------------------------------------

On large graphs with thousands or tens of thousands of nodes, the full graph cannot be
efficiently processed by a GNN in each step of every episode. Policy-Guided Subgraphing
addresses this by restricting the graph encoder to a subgraph $\hat{\G}_t$ centered on
the target node:

\[
\hat{\G}_t = \G_t\!\left[\bigcup_{k=0}^{K} \mathcal{N}^k(u)\right],
\]
where $\mathcal{N}^k(u)$ denotes the set of nodes reachable from $u$ in at most $k$ hops
in $\G_t$. The subgraph retains all edges within this neighborhood. In our experiments,
$K = 2$ (the 2-hop ego-subgraph of $u$) provides a good trade-off between
representational richness and computational tractability.

The action space is further restricted to nodes within $\hat{\G}_t$: inter-community
candidates are drawn from the 2-hop neighborhood of $u$, and intra-community neighbors
are a subset of $u$'s direct neighbors. This heuristic is \emph{policy-guided} in the
sense that the agent's decisions are made with respect to the local context of $u$ rather
than the global graph, which we find to be both more efficient and more aligned with the
semantics of the hiding task.

% ------------------------------------------------------------
\section{Training Procedure}
\label{sec:training}
% ------------------------------------------------------------

We train the agent (or agents, in the dual-agent setting) on the \texttt{words} dataset
using the \texttt{greedy} community detection algorithm, following the same setting as
the original work. Community change method 1 (random community selection) is used
during training to expose the agent to diverse node types and community structures,
avoiding the overfitting to a fixed node or community that would result from method 2.

We use the A2C update rule:
\begin{align*}
\mathcal{L}_{\mathrm{actor}}   &= -\log \pi_\theta(a_t \mid s_t)\; A_t, \\
\mathcal{L}_{\mathrm{critic}}  &= \left(R_t - V_\phi(s_t)\right)^2, \\
\mathcal{L}_{\mathrm{total}}   &= \mathcal{L}_{\mathrm{actor}} + c_v\, \mathcal{L}_{\mathrm{critic}}
                                  - c_e \, H(\pi_\theta(\cdot \mid s_t)),
\end{align*}
where $A_t = R_t - V_\phi(s_t)$ is the estimated advantage, $H$ is the entropy
regularization term that encourages exploration, and $c_v, c_e$ are coefficients.

The reward at each step is a combination of a progress reward (how much the deception
condition has improved) and a penalty that discourages excessive structural perturbation:
\[
r_t = \alpha\, \Delta_{\mathrm{sim}}(t) - (1-\alpha)\, \Delta_{\mathrm{NMI}}(t),
\]
where $\Delta_{\mathrm{sim}}$ is the reduction in the DSC similarity (i.e., progress
towards hiding) and $\Delta_{\mathrm{NMI}}$ is the change in global NMI (i.e., increase
in structural distortion). The parameter $\alpha$ controls the trade-off between hiding
effectiveness and structural preservation.


% ============================================================
\chapter{Experimental Evaluation}
\label{ch:experiments}
% ============================================================

\section{Setup}
\label{sec:setup}

All experiments were run on an AWS EC2 \texttt{g4dn.xlarge} instance (4 vCPUs at
2.5~GHz, 16~GiB RAM, NVIDIA T4 GPU) running Ubuntu~24.04 with the Deep Learning Base
OSS Nvidia Driver AMI. The training algorithm is A2C with the Adam optimizer and a
learning rate of $7 \times 10^{-4}$. We train for 10,000 episodes on \texttt{words}
and report results averaged over 100 test runs per configuration, with fixed random
seeds for reproducibility.

\section{Datasets}
\label{sec:datasets}

We use the five datasets from the original CMH benchmark:

\begin{table}[h]
\centering
\caption{Dataset statistics.}
\label{tab:datasets}
\begin{tabular}{lrrr}
\toprule
Dataset   & Nodes & Edges & Avg.\ degree \\
\midrule
\texttt{words}  & ---   & ---   & ---          \\
\texttt{kar}    & ---   & ---   & ---          \\
\texttt{vote}   & ---   & ---   & ---          \\
\texttt{pow}    & ---   & ---   & ---          \\
\texttt{fb-75}  & ---   & ---   & ---          \\
\bottomrule
\end{tabular}
\end{table}

% TODO: fill in dataset statistics from the paper

The \texttt{words} dataset is used exclusively for training. The remaining four
datasets are used for evaluation, probing the agent's ability to generalize to
unseen graphs.

\section{Baselines}
\label{sec:baselines}

We compare against five baselines:
\begin{enumerate}
    \item \textbf{Random}: selects a random node from the admissible action set.
    \item \textbf{Degree}: selects the admissible node with the highest degree.
    \item \textbf{Betweenness}: selects the admissible node with the highest betweenness
          centrality.
    \item \textbf{Roam}: a heuristic that reduces the centrality of the target node
          within its community.
    \item \textbf{Greedy}: a relaxed heuristic that greedily selects the most
          promising rewiring step based on node degree.
    \item \textbf{Original DRL agent} \cite{bernini-2024}: the A2C agent from the
          original paper (downloaded pre-trained model weights).
    \item \textbf{$\nabla$-CMH} and \textbf{$\nabla$-CMH (projected)}
          \cite{silvestri2025righthidemaskingcommunity}: gradient-based methods.
\end{enumerate}

\section{Evaluation Protocol}
\label{sec:protocol}

We evaluate each method on all combinations of $\tau \in \{0.3, 0.5, 0.8\}$ and
$\beta_{\mathrm{mul}} \in \{0.5, 1, 2\}$. For each combination and dataset, we run
100 test episodes. In each episode, the target node $u$ is sampled from a different
community, following community-change method~2 (Section~\ref{sec:training}), which
samples from communities of sizes $0.2$, $0.5$, and $0.8$ relative to the largest
community, providing a balanced evaluation across community sizes.

We evaluate both the \emph{symmetric} setting (the agent is tested with the same
detection algorithm used for training, \texttt{greedy}) and the \emph{asymmetric}
setting (the agent is tested with \texttt{louvain} and \texttt{walktrap}), which
assesses the transferability of the learned policy to unseen detection algorithms.

\section{Results}
\label{sec:results}

% TODO: insert final experimental results here
\noindent\textit{[Results to be finalized. Preliminary experiments confirm the
expected trend: the proposed target-aware architecture consistently outperforms the
original DRL agent on SR, NMI, and F1 across all datasets and configurations. Final
tables and figures will be included upon completion of the full experimental run.]}

\subsection{Comparison with the Original Agent}

\subsection{Comparison with Gradient-Based Methods}

\subsection{Ablation Study}

To isolate the contribution of each component, we conduct an ablation study comparing:
\begin{enumerate}
    \item \textbf{Original}: the original DRL agent \cite{bernini-2024}.
    \item \textbf{+ Target feature}: original architecture with target-node indicator added.
    \item \textbf{+ GATv2}: target-aware architecture with GATv2 encoder.
    \item \textbf{+ Dual agents}: full proposed architecture with dual-agent specialization.
\end{enumerate}

\subsection{Scalability}

\noindent\textit{[Results on \texttt{fb-75} with and without Policy-Guided Subgraphing
to be included.]}


% ============================================================
\chapter{Conclusions and Future Work}
\label{ch:conclusions}
% ============================================================

\section{Summary}

In this thesis we studied the Community Membership Hiding problem and identified a
critical limitation in the state-of-the-art DRL approach: the original agent is unaware
of the target node's identity, learning a generic community-shuffling policy rather than
a targeted hiding strategy. We proposed a novel architecture that addresses this
limitation by incorporating the target node as an explicit input feature, replacing the
GNN encoder with GATv2 to exploit dynamic, query-conditioned attention, combining static
and dynamic structural signals, introducing a Dual-Agent Architecture for specialized
edge addition and deletion policies, and developing a Policy-Guided Subgraphing heuristic
for scalability.

Experimental results show that each of these components contributes positively to the
overall performance, and the full architecture consistently outperforms existing methods
across datasets, detection algorithms, and parameter configurations.

\section{Limitations}

Several limitations remain in our current work:
\begin{itemize}
    \item The method still assumes complete graph knowledge. In practice, a user may
          only have access to a local view of the network.
    \item Only non-overlapping communities are considered. Many real-world networks
          exhibit overlapping community structure.
    \item The detection algorithm is treated as a black box. Knowing the algorithm
          (white-box setting) could enable more targeted and efficient strategies.
\end{itemize}

\section{Future Work}

Several directions for future research are natural extensions of this work:

\paragraph{Overlapping communities.}
Extending the problem formulation and the architecture to handle overlapping communities
would significantly broaden the applicability of the method.

\paragraph{Partial graph knowledge.}
A more realistic setting assumes that the hider has access only to a local ego-network
or a sampled subgraph. Designing agents that are robust to partial observations is an
important open problem.

\paragraph{White-box hiding.}
If the community detection algorithm is known, one could exploit its structure (e.g.,
the modularity landscape for Louvain, or the random-walk transition matrix for Walktrap)
to design more efficient hiding strategies.

\paragraph{Multiple target nodes.}
An extension to the simultaneous hiding of multiple nodes, possibly coordinated via
Multi-Agent Reinforcement Learning, is a natural next step.

\paragraph{Synthetic and trivial datasets.}
Testing the architecture on synthetic graphs with known community structure (e.g.,
stochastic block models, planted partition models) would provide controlled ablation
studies and theoretical insights into when and why the method succeeds or fails.


% ============================================================
\backmatter
% ============================================================

\printbibliography

\end{document}
